% ----------------------------------------------------------------
% AMS-LaTeX definitions for a Dissertation/Thesis ****************
% Template prepared by T.A. McWalter (March 2009)
% ----------------------------------------------------------------
\documentclass[a4paper,11pt,oneside]{report}

% PACKAGES -------------------------------------------------------

\usepackage{amsmath}
\usepackage{amsthm}
\usepackage{amsfonts}
\usepackage{amssymb}
\usepackage{fancyhdr}
\usepackage{footnote}
\usepackage{ifpdf}
\usepackage{mathrsfs}
\usepackage[numbers]{natbib}
\usepackage{xspace}

% GRAPHICS HANDLING ----------------------------------------------

\ifpdf
    \usepackage[pdftex]{graphicx}
    \DeclareGraphicsExtensions{.pdf, .jpg, .tif}
    \newcommand{\bb}{viewport}
\else
    \usepackage{graphicx}
    \DeclareGraphicsExtensions{.eps, .jpg}
    \newcommand{\bb}{bb}
\fi

% OTHER SET-UP ---------------------------------------------------

\pagestyle{fancy}
\lhead{\rightmark}
\chead{}
\rhead{\thepage}
\cfoot{}

\setlength{\topmargin}{0.5cm}
\setlength{\headsep}{1cm}
\setlength{\textwidth}{14.0cm}
\setlength{\headwidth}{14.0cm}
\setlength{\headheight}{14pt}
\setlength{\textheight}{22cm}
\setlength{\oddsidemargin}{1.3cm}
\setlength{\evensidemargin}{0cm}

\newcommand{\setlinespacing}[1]
           {\renewcommand{\baselinestretch}{#1}\small\normalsize}

% THEOREM STYLES -------------------------------------------------

\theoremstyle{plain}
\newtheorem{thm}{Theorem}[chapter]
\newtheorem{prop}[thm]{Proposition}
\newtheorem{cor}[thm]{Corollary}
\newtheorem{con}[thm]{Conjecture}
\newtheorem{lem}[thm]{Lemma}
\theoremstyle{definition}
\newtheorem{defn}[thm]{Definition}
\newtheorem{ass}[thm]{Assumption}
\newtheorem{exam}[thm]{Example}

% MACRO DEFINITIONS ----------------------------------------------

\def\n{\nonumber\\}

% Statistical
\def\E[#1]{\Tilde{\mathbb{E}}\!\left[#1\right]}
\def\Et[#1]{\Tilde{\mathbb{E}}_t\!\left[#1\right]}
\def\Var[#1]{\Tilde{\mathbb{V}}\mathtt{ar}\!\left[#1\right]}
\def\Prob(#1){\Tilde{\mathbb{P}}\!\left(#1\right)}

% Derivatives
\def\S{\ensuremath{S_t}\xspace}
\def\dP{\ensuremath{dP_t}\xspace}
\def\dS{\ensuremath{dS_t}\xspace}
\def\dt{\ensuremath{dt}\xspace}
\def\dW{\ensuremath{dW_t}\xspace}
\def\dq{\ensuremath{dq_t}\xspace}
\def\Vt{\frac{\partial V}{\partial t}}
\def\Vs{\frac{\partial V}{\partial S}}
\def\Vss{\frac{\partial^2 V}{\partial S^2}}

% Differences
\def\DS{\ensuremath{\delta S_t}\xspace}
\def\Dt{\ensuremath{\delta t}\xspace}
\def\DW{\ensuremath{\delta W_t}\xspace}
\def\Dq{\ensuremath{\delta q_t}\xspace}

% Finite difference
\def\Vjn{\ensuremath{V^{n}_j}\xspace}
\def\Vjnm{\ensuremath{V^{n - 1}_j}\xspace}
\def\Vjnp{\ensuremath{V^{n + 1}_j}\xspace}
\def\Vjmn{\ensuremath{V^{n}_{j - 1}}\xspace}
\def\Vjpn{\ensuremath{V^{n}_{j + 1}}\xspace}
\def\Vjmnm{\ensuremath{V^{n - 1}_{j - 1}}\xspace}
\def\Vjpnm{\ensuremath{V^{n - 1}_{j + 1}}\xspace}
\def\Vjmnp{\ensuremath{V^{n + 1}_{j - 1}}\xspace}
\def\Vjpnp{\ensuremath{V^{n + 1}_{j + 1}}\xspace}

\def\Von{\ensuremath{V^{n}_{0}\xspace}}
\def\Vomn{\ensuremath{V^{n}_{-1}\xspace}}
\def\Vopn{\ensuremath{V^{n}_{1}\xspace}}
\def\VJn{\ensuremath{V^{n}_{J}\xspace}}
\def\VJmn{\ensuremath{V^{n}_{J - 1}\xspace}}
\def\VJpn{\ensuremath{V^{n}_{J + 1}\xspace}}
\def\Vonm{\ensuremath{V^{n - 1}_{0}\xspace}}
\def\Vomnm{\ensuremath{V^{n - 1}_{-1}\xspace}}
\def\Vopnm{\ensuremath{V^{n - 1}_{1}\xspace}}
\def\VJnm{\ensuremath{V^{n - 1}_{J}\xspace}}
\def\VJmnm{\ensuremath{V^{n - 1}_{J - 1}\xspace}}
\def\VJpnm{\ensuremath{V^{n - 1}_{J + 1}\xspace}}

\def\Sj{\ensuremath{S_j}\xspace}
\def\Xj{\ensuremath{X_j}\xspace}

% Binomial model
\def\u{\ensuremath{u}\xspace}
\def\d{\ensuremath{d}\xspace}
\def\o{\ensuremath{o}\xspace}

\def\pu{\ensuremath{p_\u}\xspace}
\def\pd{\ensuremath{p_\d}\xspace}
\def\po{\ensuremath{p_\o}\xspace}

% Binomial model valuation
\def\V{\ensuremath{V_t}\xspace}
\def\Vu{\ensuremath{V^{u}_t}\xspace}
\def\Vd{\ensuremath{V^{d}_t}\xspace}
\def\X{\ensuremath{X_t}\xspace}
\def\P{\ensuremath{\Pi_t}\xspace}
\def\D{\ensuremath{\Delta_t}\xspace}

% Sundry
\def\th{\textsuperscript{th}\xspace}
\def\half{\tfrac{1}{2}}

% PREAMBLE -------------------------------------------------------

\title{Valuation of Convertible Bonds with Credit Risk}

\author{David Naylor\\[2cm]
Programme in Advanced Mathematics of Finance,\\
University of the Witwatersrand,\\
Johannesburg.\\[2cm]
\includegraphics[width=4cm]{../common/WitsColourLogo}\\[2cm]
}

% MAIN DOCUMENT --------------------------------------------------

\begin{document}

\maketitle \setlinespacing{1.20}

% DECLARATION ----------------------------------------------------

\def\abstractname{\Huge Declaration}
\begin{abstract}
I declare that this project is my own, unaided work. It is being submitted as partial fulfilment of the Degree of Bachelor of Science with Honours in the University of the Witwatersrand, Johannesburg. It has not been submitted before for any degree or examination in any other University.\\[3cm]
\noindent\rule{5cm}{0.5pt}\\[1cm]
\noindent \today
\end{abstract}

% ABSTRACT -------------------------------------------------------

\def\abstractname{\Huge Abstract}
\begin{abstract}
A convertible bond is a complex derivative that cannot be priced as a simple combination of bond and stock components.  Convertible bonds can be broken down as a bond with two embedded options (a put option for the investor and a call option for the issuer) and an option to convert the bond into stock.  Due to the multiple continuous options, the pricing of the convertible bond is path dependent.

This research project explores and implements a binary tree and finite difference scheme to price the convertible bond, taking into account credit risk.
\end{abstract}

% ACKNOWLEDGEMENTS -----------------------------------------------

\def\abstractname{\Huge Acknowledgements}
\begin{abstract}
Place acknowledgements here.
\end{abstract}

% TABLE OF CONTENTS ----------------------------------------------

\setlinespacing{1}
\pagenumbering{roman}
\tableofcontents
\listoffigures
\listoftables
\clearpage
\setlinespacing{1.20}
\pagenumbering{arabic}
\setcounter{page}{1}

% CONTENT --------------------------------------------------------

\chapter{Introduction}
What this paper shows and does.

\section{Convertible Bonds with Credit Risk}
Describe and define (in maths) what a convertible bond with credit risk is.

\section{Literature Review}
Describe the current state of literature w/ Convertible Bonds with Credit Risk,
possibly split this and include it at the beginning of each relevant section?


\chapter{Finite Difference Model}
Describe the finite difference model based on the paper by \citet{AFV03}.

\section{Definition}
Consider a stock with the following stochastic model for stock price movements:
\begin{equation}
 \dS = \mu\S\dt + \sigma\S\dW - \eta\S\dq
\end{equation}
If we impose that the risk neutral expectation of the stock price evolve at the risk free rate $\E[\dS] = r\S\dt$ then $\mu = (r + \lambda\eta)$ and the differential equation becomes:
\begin{equation}
 \dS = (r + \lambda\eta)\S\dt + \sigma\S\dW - \eta\S\dq
\end{equation}
where:

\begin{tabular}{ll}
 $\mu$          & is the drift rate of the stock price \\
 $r$            & is the risk free rate \\
 $\sigma$       & is the log-volatility of the stock price \\
 $\lambda$      & is the hazard rate \\
 $\eta$         & is the percentage drop in stock price on a default event \\
 \dW            & is a Wiener process \\
 \dq            & is a Poisson jump process where the first jump is the default event \\
\end{tabular}

\section{Derivation}
Consider a portfolio, \P, that has a long position in the derivative, \V, and a short position of \D in the stock, \S.  The residual $\P = \V - \D\S$ is invested at the risk neutral rate.  This yields a differential equation for \P:
\begin{align}
 \dP &= r\P\dt \n
     &= r(\V - \D\S)\dt \label{eq:rPdt}
\end{align}

In the event of default the portfolio holdings $\P^x$ have a long position in the derivative with residual value \X and a short position of \D in the residual stock with value of $(1 - \eta)\S$.  Thus the difference of the portfolio holdings when going into default is:
\begin{align}
 \dP^x &= \P^x - \P \n
        &= \X - (\V - \eta\D\S) \label{eq:dPx}
\end{align}

At any given instant the portfolio can either default or evolve in the normal stochastic manner.  The probability of default in any given instant is $\lambda\dt$ and the probability of non-default is $(1 - \lambda)\dt$.  Taking the expectation of the differential of \P and applying Ito's lemma to the stochastic component:
\begin{align}
 \Et[\dP] &= (1 - \lambda\dt)[\left(\Vt + \mu\S\Vs + \half\sigma^2\S^2\Vss - \mu\S\D\right)\dt + \n
          &\qquad \left(\sigma\S\Vs - \sigma\S\D\right)\dW] + \lambda\dt(\X - (\V - \eta\D\S))
\end{align}
By choosing \D to eliminate the \dW term, $\D = \Vs$, the differential simplifies to:
\begin{equation}
 \Et[\dP] = \left(\Vt + \half\sigma^2\S^2\Vss\right)\dt + \lambda\dt(\X - (\V - \eta\D\S)) \label{eq:EdP}
\end{equation}

Equating the differential \eqref{eq:rPdt} and the expected differential \eqref{eq:EdP} for \P the partial differential equation for a derivative with credit risk is derived:
\begin{align}
 r(\V - \D\S)\dt &= \left(\Vt + \half\sigma^2\S^2\Vss\right)\dt + \lambda\dt(\X - (\V - \eta\D\S)) \n
 (r + \lambda)\V &= \Vt + (r + \lambda\eta)\S\Vs + \half\sigma^2\S^2\Vss + \lambda\X
\end{align}

\section{Finite Difference}
The derived PDE can be solved numerically by constructing a lattice of stock prices and time.  \Vjn will denote the lattice value at the j\th stock node and n\th time node where the step between stock nodes is \DS and the step between time nodes is \Dt.  \Sj will denote the stock value at the j\th stock node and \Xj will denote the defaulted value at the j\th stock node.

All finite difference equations derived here can be expressed as functions of $\alpha$, $\beta$ and $\kappa$ where:
\begin{align*}
 \alpha &= \half\sigma^2\Sj^2\frac{\Dt}{\DS^2} - (r + \lambda\eta)\Sj\frac{\Dt}{2\DS} \\
 \beta  &= - (r + \lambda) - \sigma^2\Sj^2\frac{\Dt}{\DS^2} \\
 \kappa &= \half\sigma^2\Sj^2\frac{\Dt}{\DS^2} + (r + \lambda\eta)\Sj\frac{\Dt}{2\DS} \\
\end{align*}
and it is assumed that for boundary nodes:
\begin{align*}
 \frac{\Von - \Vomn}{\DS} &= \frac{\Vopn - \Von}{\DS} \\
                    \Vomn &= 2\Von - \Vopn \\
 \frac{\VJn - \VJmn}{\DS} &= \frac{\VJpn - \VJn}{\DS} \\
                    \VJpn &= 2\VJn - \VJmn \\
\end{align*}

\subsection{Explicit Scheme}
Using the backwards difference equation for the temporal derivative the finite difference equation is:
\begin{displaymath}
 (r + \lambda)\Vjn = \frac{\Vjn - \Vjnm}{\Dt} + (r + \lambda\eta)\Sj\frac{\Vjpn - \Vjmn}{2\DS} + \half\sigma^2\Sj^2\frac{\Vjmn - 2\Vjn + \Vjpn}{\DS^2} + \lambda\Xj
\end{displaymath}
\begin{align}
 \Vjnm &= \alpha\Vjmn + (1 + \beta)\Vjn + \kappa\Vjpn + \lambda\Dt\Xj \\
 \Vonm &= (1 + \beta + 2\alpha)\Von + (\kappa - \alpha)\Vopn + \lambda\Dt\Xj \\
 \VJnm &= (\alpha - \kappa)\Vjmn + (1 + \beta - 2\kappa)\VJn + \lambda\Dt\Xj
\end{align}

\subsection{Implicit Scheme}
Using the backwards difference equation for the temporal derivative and evaluating the spacial differentials at the previous time step the finite difference equation is:
\small\begin{displaymath}
 (r + \lambda)\Vjnm = \frac{\Vjn - \Vjnm}{\Dt} + (r + \lambda\eta)\Sj\frac{\Vjpnm - \Vjmnm}{2\DS} + \half\sigma^2\Sj^2\frac{\Vjmnm - 2\Vjnm + \Vjpnm}{\DS^2} + \lambda\Xj
\end{displaymath}\normalsize
\begin{align}
 -\alpha\Vjmnm + (1 - \beta)\Vjnm - \kappa\Vjpnm &= \Vjn + \lambda\Dt\X \\
 (1 - \beta - 2\alpha)\Vonm - (\kappa - \alpha)\Vopnm &= \Von + \lambda\Dt\X \\
 -(\alpha - \kappa)\VJmnm + (1 - \beta - 2\kappa)\VJnm &= \VJn + \lambda\Dt\X
\end{align}

\subsection{Crank-Nicolson Scheme}
Using the backwards difference equation for the temporal derivative and taking the average of the spacial differentials at the current and previous time step the finite difference equation is:
\small\begin{align*}
 (r + \lambda)\half(\Vjn + \Vjnm) &= \frac{\Vjn - \Vjnm}{\Dt} + (r + \lambda\eta)\Sj\frac{1}{2}\left(\frac{\Vjpn - \Vjmn}{2\DS} + \frac{\Vjpnm - \Vjmnm}{2\DS}\right) + \\
  &\quad \half\sigma^2\Sj^2\frac{1}{2}\left(\frac{\Vjmn - 2\Vjn + \Vjpn}{\DS^2} + \frac{\Vjmnm - 2\Vjnm + \Vjpnm}{\DS^2}\right) + \lambda\Xj
\end{align*}
\begin{align}
 -\alpha\Vjmnm + (2 - \beta)\Vjnm - \kappa\Vjpnm &= \alpha\Vjmn + (2 + \beta)\Vjn + \kappa\Vjpn + 2\lambda\Dt\X \\
 (2 - \beta - 2\alpha)\Vonm - (\kappa - \alpha)\Vopnm &= (2 + \beta + 2\alpha)\Von + (\kappa + \alpha)\Vopn + 2\lambda\Dt\X \\
 -(\alpha - \kappa)\VJmnm + (2 - \beta - 2\kappa)\VJnm &= (\alpha - \kappa)\VJmn + (2 + \beta + 2\kappa)\VJn + 2\lambda\Dt\X
\end{align}\normalsize


\chapter{Binomial Model with Credit Risk}
The Binomial Model with Credit Risk was first derived by \citet{MK12}, who showed that this model converges, in continuous time, to the \citet{AFV03} model.  \Citet{MK12} also showed the valuation method of this model is the same as in the classical binomial model.

TODO: add further explanation and references.

\section{Definition}
Consider a binomial model that has time step \Dt, up and down steps, \u and \d with probability \pu and \pd\footnote{$\pd = 1 - \pu - \po = e^{-\lambda\Dt} - \pu$} respectively, and where the probability of default is \po\footnote{The time till the first jump follows an exponential distribution with intensity $\lambda$ and has probability $\Prob(q_{t + \dt} = 1 | q_t = 0) = \po = 1 - e^{-\lambda\Dt}$}.  The mean and variance of this model is: (TODO: insert picture illustration)
\begin{align}
 \E[\frac{S_{t + \Dt}}{\S}]     &= \u\pu + \d\pd + (1 - \eta)\po \n
                                &= \u\pu + \d(e^{-\lambda\Dt} - \pu) + (1 - \eta)(1 - e^{-\lambda\Dt}) \label{eq:BM_E} \\
 \Var[\frac{S_{t + \Dt}}{\S}]   &= \u^2\pu + \d^2\pd + (1 - \eta)^2\po - \E[\frac{\DS}{\S}]^2 \n
                                &= \u^2\pu + \d^2(e^{-\lambda\Dt} - \pu) + (1 - \eta)^2(1 - e^{-\lambda\Dt}) - \E[\frac{S_{t+\Dt}}{\S}]^2 \label{eq:BM_Var}
\end{align}

The mean and variance of the stock's stochastic differential equation is:
\begin{align}
 \E[\frac{\dS}{\S}]     &= r\dt \\
 \E[\frac{S_{t + \dt}}{\S}] &= \E[\frac{\S}{\S} + \frac{\dS}{\S}] \n
                        &= 1 + \E[\frac{\dS}{\S}] \label{eq:SM_E} \\
 \Var[\frac{\dS}{\S}]   &= (\sigma^2 + \lambda\eta^2)\dt \\
 \Var[\frac{S_{t + \dt}}{\S}] &= \Var[\frac{\S}{\S} + \frac{\dS}{\S}] \n
                        &= \Var[\frac{\dS}{\S}] \label{eq:SM_Var}
\end{align}


\section{Derivation}
If one equates the first moment of the stochastic model \eqref{eq:SM_E}\footnote{$\E[\frac{S_{t + \Dt}}{\S}] = e^{r\Dt}$}, in \Dt time, with that of the binomial model \eqref{eq:BM_E} then:
\begin{align}
 e^{r\Dt}       &= \u\pu + \d(e^{-\lambda\Dt} - \pu) + (1 - \eta)(1 - e^{-\lambda\Dt}) \n
 \pu(\u - \d)   &= e^{r\Dt} - \d e^{-\lambda\Dt} - (1 - \eta)(1 - e^{-\lambda\Dt}) \label{eq:p(u-d)}\\
 \pu            &= \frac{e^{r\Dt} - \d e^{-\lambda\Dt} - (1 - \eta)(1 - e^{-\lambda\Dt})}{\u - \d} \\
 \therefore\quad\pd &= -\frac{e^{r\Dt} - \u e^{-\lambda\Dt} - (1 - \eta)(1 - e^{-\lambda\Dt})}{\u - \d}
\end{align}
and if one equate the second moment about the mean of the stochastic model \eqref{eq:SM_Var}, in \Dt time, with that of the binomial model \eqref{eq:BM_Var} then\footnote{Substituting equation \eqref{eq:p(u-d)} into \eqref{eq:Var=1} one gets \eqref{eq:Var=2}}:
\begin{align}
 (\sigma^2 + \lambda\eta^2)\Dt &= \u^2\pu + \d^2(e^{-\lambda\Dt} - \pu) + (1 - \eta)^2(1 - e^{-\lambda\Dt}) - \E[\frac{S_{t + \Dt}}{\S}]^2 \n
                        &= (\u^2 - \d^2)\pu + \d^2 e^{-\lambda\Dt} +  (1 - \eta)^2(1 - e^{-\lambda\Dt}) - e^{-2r\Dt} \label{eq:Var=1} \\
                        &= (\u + \d)(e^{r\Dt} - \d e^{-\lambda\Dt} - (1 - \eta)(1 - e^{-\lambda\Dt})) \label{eq:Var=2} \\
                        &\qquad + \d^2 e^{-\lambda\Dt} + (1 - \eta)^2(1 - e^{-\lambda\Dt}) - e^{-2r\Dt} \n
                        &= (\u + \d)(e^{r\Dt} - (1 - \eta)(1 - e^{-\lambda\Dt})) - \u\d e^{-\lambda\Dt} \label{eq:Var=} \\
                        &\qquad + (1 - \eta)^2(1 - e^{-\lambda\Dt}) - e^{-2r\Dt} \nonumber
\end{align}
If one assumes $\Dt^2 = 0$, $\u\d = 1$ and $\u = e^{\sqrt{A\Dt}}$ and the Taylor series expansion is taken for all exponential terms, then:
\begin{align}
 \u                     &= 1 + \sqrt{A\Dt} + \frac{A\Dt}{2!} + \frac{(A\Dt)^{\frac{3}{2}}}{3!} \n
 \d                     &= 1 - \sqrt{A\Dt} + \frac{A\Dt}{2!} - \frac{(A\Dt)^{\frac{3}{2}}}{3!} \n
 \u + \d                &= 2 + A\Dt \label{eq:u+d}
\end{align}
and substituting \eqref{eq:u+d} into \eqref{eq:Var=} with the appropriate expansions then:
\begin{align}
 (\sigma^2 + \lambda\eta^2)\Dt &= (2 + A\Dt)(1 + r\Dt - \lambda\Dt(1 - \eta)) - (1 - \lambda\Dt) \n
                        &\qquad + \lambda\Dt(1 - \eta)^2 - (1 - 2r\Dt) \n
                        &= A\Dt + \lambda\Dt - \lambda\Dt(1 - \eta)(1 + \eta) \n
                        &= A\Dt + \lambda\eta^2\Dt \n
 A\Dt                   &= \sigma^2\Dt
\end{align}

\section{Limits}
The time step \Dt is assumed to be strictly positive ($\Dt > 0$), as time is monotonically increasing.

Although $\sigma$ could be positive or negative, the formul\ae\ use $\sigma$ where, strictly, $|\sigma|$ should be used.  If $\sigma = 0$ then the model is no longer useful.  The restriction $\sigma > 0$ may be introduced for convenience and without loss of generality.

An implicit assumption this model makes is that \pu, \pd and \po are a valid probability measure.  This imposes the following conditions:
\begin{align}
 \min(\pu, \pd, \po) &\ge 0 \\
 \pu + \pd + \po &= 1 \label{eq:bin_prob_limit}
\end{align}

Consider \po, the probability of default.  Its value is equal to the cumulative density function of an exponential distribution with intensity $\lambda$.  As such \po is already confined to the interval $[0, 1]$ with the restriction that $\lambda > 0$.  If $\lambda = 0$, no default is possible, then \po will have a value of 0 and is a valid probability thus the restriction on $\lambda$ can be relaxed to $\lambda \ge 0$.

Consider \pd, the probability of a down movement.  Its value is defined in terms of \pu and \po.  If \pu and \po are valid probabilities, within bounds, then so will \pd.

Consider \pu, the probability of an up movement.  Taking into account the definition of \pd, this requires that \pu satisfies:
\begin{align}
 0 &\le \pu \le e^{-\lambda \Dt} \\
 0 &\le \frac{e^{r\Dt} - \d e^{-\lambda\Dt} - (1 - \eta)(1 - e^{-\lambda\Dt})}{\u - \d} \le e^{-\lambda \Dt} \n
 0 &\le e^{r\Dt} - \d e^{-\lambda\Dt} - (1 - \eta)(1 - e^{-\lambda\Dt}) \le e^{-\lambda \Dt}(\u - \d) \n
 e^{-\lambda\Dt}(\d - (1 - \eta)) &\le e^{r\Dt} - (1 - \eta) \le e^{-\lambda\Dt}(\u - (1 - \eta)) \n
 \ln\left(\frac{d - (1 - \eta)}{e^{r\Dt} - (1 - \eta)}\right) &\le \lambda\Dt \le \ln\left(\frac{u - (1 - \eta)}{e^{r\Dt} - (1 - \eta)}\right) \\
 \Dt &\le \frac{1}{\lambda}\ln\left(\frac{u - (1 - \eta)}{e^{r\Dt} - (1 - \eta)}\right)
\end{align}
Using the limit on both \Dt and $\lambda$ the inequalities around the $\lambda\Dt$ are also limited thus:
\begin{align}
 0 &\le \ln\left(\frac{u - (1 - \eta)}{e^{r\Dt} - (1 - \eta)}\right) \n
 e^{r\Dt} - (1 - \eta) &\le u - (1 - \eta) \n
 r\Dt &\le \sigma \sqrt{\Dt} \n
 \Dt &\le \frac{\sigma^2}{r^2}
\end{align}
The inequality of $\ln\left(\frac{d - (1 - \eta)}{e^{r\Dt} - (1 - \eta)}\right) \le 0$ results in the same inequality as above.

Based on the probability limits \eqref{eq:bin_prob_limit} no restriction is placed on $r$, the risk free rate, or $\eta$, the percentage drop in stock price.

\section{Valuation}
Consider a portfolio, \P, that has a long position in the derivative, \V, and a short position of \D in the stock, \S.  The residual $\P = \V - \D\S$ is invested at the risk neutral rate.  At time $t + \Dt$ one should have:
\begin{equation*}
 \Pi_{t + \Dt} =
 \begin{cases}
  \Vu - \D\S u          & \text{with probability } \pu \\
  \Vd - \D\S d          & \text{with probability } \pd\\
  \X - (1 - \eta)\S     & \text{with probability } \po\\
 \end{cases}
\end{equation*}
where

\begin{tabular}{ll}
 $\Vu$          & is the value of the portfolio on an up-step at time $t + \Dt$ \\
 $\Vd$          & is the value of the portfolio on a down-step at time $t + \Dt$ \\
 $\X$           & is the value of default at time $t$ \\
\end{tabular}

If one wishes to hedge against up and down movements of the stock then:
\begin{align}
    \Vu - \D\S u &= \Vd - \D\S d \n
     \D\S(d - u) &= \Vd - \Vu \n
              \D &= \frac{\Vu - \Vd}{\S(u - d)}
\end{align}

The required value of \P at time $t + \Dt$ is:
\begin{align}
 \Pi_{t + \Dt}  &= \P e^{r\Dt} \n
                &= \left(\V - \frac{\Vu - \Vd}{u - d}\right) e^{r\Dt} \label{eq:P}
\end{align}
as \P is invested at the risk free rate, and taking the expected value of $\Pi_{t + \Dt}$, which must equal \eqref{eq:P}, one gets:
\begin{align}
 \left(\V - \frac{\Vu - \Vd}{u - d}\right) e^{r\Dt} &= \Et[\Pi_{t + \Dt}] \n
        &= \left(\Vu - \frac{\Vu - \Vd}{u - d}u\right) e^{-\lambda\Dt} \n
        &\qquad + \left(\X - \frac{\Vu - \Vd}{u - d}(1 - \eta)\right)(1 - e^{-\lambda\Dt}) \n
        &= \Vu\left(e^{-\lambda\Dt} - \frac{ue^{-\lambda\Dt} + (1 - \eta)(1 - e^{-\lambda\Dt})}{u - d}\right) \n
        &\qquad + \Vd\left(\frac{ue^{-\lambda\Dt} + (1 - \eta)(1 - e^{-\lambda\Dt})}{u - d}\right) + \X(1 - e^{-\lambda\Dt}) \n
 \V e^{r\Dt} &= \Vu\frac{e^{r\Dt} - de^{-\lambda\Dt} - (1 - \eta)(1 - e^{-\lambda\Dt})}{u - d} \n
        &\qquad + \Vd\frac{ue^{-\lambda\Dt} + (1 - \eta)(1 - e^{-\lambda\Dt}) - e^{r\Dt}}{u - d} + \X(1 - e^{-\lambda\Dt}) \n
 \V     &= e^{-r\Dt}(\Vu\pu + \Vd\pd + \X\po)
\end{align}
thus the pricing method of this binomial model is the same as the pricing method of the classical binomial model:
\begin{align}
 V_t    &= e^{-r\Dt}(\Vu\pu + \Vd\pd + \X\po) \n
        &= e^{-r\Dt}\Et[V_{t + \Dt}]
\end{align}

\section{Valuation with Coupon Cash-flow}
Consider a portfolio, as above, with the add of coupon cash flows.  At time $t + \Dt$ one should have:
\begin{equation*}
 \Pi_{t + \Dt} =
 \begin{cases}
  \Vu - \D\S u + c_i e^{t + \Dt - t^{c}_i}      & \text{with probability } \pu \\
  \Vd - \D\S d + c_i e^{t + \Dt - t^{c}_i}      & \text{with probability } \pd\\
  \X - (1 - \eta)\S + c_i (1 - q_{t^{c}_i})e^{t + \Dt - t^{c}_i} & \text{with probability } \po\\
 \end{cases}
\end{equation*}
where

\begin{savenotes}
\begin{tabular}{ll}
 $c_i$          & is the i\textsuperscript{th} coupon value \\
 $t^{c}_i$      & is the time of the i\textsuperscript{th} coupon, with $t^{c}_i \in (t, t + \Dt)$ \\
 $q_{t^{c}_i}$  & is an indicator variable that default happens after coupon payment\footnote{$\Prob(q_{t^{c}_i} = 0 | q_s = 1, s \in (t, t + \Dt]) = \frac{e^{-\lambda(t^{c}_i - t)} - e^{-\lambda\Dt}}{1 - e^{-\lambda\Dt}}$, the probability that the coupon arrives before default occurs, given that default occurs within the interval}
\end{tabular}
\end{savenotes}

Using the above method it is trivial to show that:
\begin{align}
 V_t    &= e^{-r\Dt}(\Vu\pu + \Vd\pd + \X\po) + c_i e^{-(r + \lambda)(t^{c}_i - t)} \n
        &= e^{-r\Dt}\Et[V_{t + \Dt}] + c_i e^{-(r + \lambda)(t^{c}_i - t)}
\end{align}
if, however, $t^{c}_i = t$, thus $c_i$ arrives with certainty at $t$, it is trivially shown that:
\begin{align}
 V_t    &= e^{-r\Dt}(\Vu\pu + \Vd\pd + \X\po) + c_i \n
        &= e^{-r\Dt}\Et[V_{t + \Dt}] + c_i
\end{align}

\section{Parameters and Formul\ae\ of Model}
Based on the above Table \ref{tab:bin_parameter} lists the parameters required for, Table \ref{tab:bin_limit} lists the limits of, and Table \ref{tab:bin_formulae} lists the formul\ae\ to price, this binomial table.

\begin{table}[ht]
\centering
\caption{Parameters of the Binomial Model with Credit Risk}
\begin{tabular}{|r|l|}\hline
 \textbf{Parameter} & \textbf{Description} \\\hline
 $r$            & Risk free rate \\\hline
 $\sigma$       & Log-volatility of the stock price \\\hline
 $\lambda$      & Hazard rate of default \\\hline
 $\eta$         & Percentage drop of stock price in a default event \\\hline
 \Dt            & Time-step \\\hline
 \Vu            & Value of portfolio for an up-step at $t + \Dt$ \\\hline
 \Vd            & Value of portfolio for a down-step at $t + \Dt$ \\\hline
 \X             & Value of defaulted portfolio at $t$ \\\hline
 $c_i$          & Value of the i\textsuperscript{th} coupon payment \\\hline
 $t^{c}_i$      & Time of the i\textsuperscript{th} coupon payment \\\hline
\end{tabular}
\label{tab:bin_parameter}
\end{table}

\begin{savenotes}
\begin{table}[ht]
\centering
\caption{Limits of the Binomial Model with Credit Risk}
\begin{tabular}{|c|l|}\hline
 \textbf{Limit} & \textbf{Description} \\\hline
 $0 < \sigma$ & Volatility must be positive \\\hline
 $0 \le \lambda$ & Hazard rate must be non-negative \\\hline
 $0 < \Dt$ & Time step must be positive \\\hline
 $\Dt \le \frac{\sigma^2}{r^2}$ & Time step must be small enough to handle volatility\footnote{If $r = 0$ then this inequality is not applicable.} \\\hline
 $\Dt \le \frac{1}{\lambda}\ln\left(\frac{u - (1 - \eta)}{e^{r\Dt} - (1 - \eta)}\right)$ & Time step must be small enough to handle hazard rate\footnote{If $\lambda = 0$ then this inequality is not applicable.} \\\hline
\end{tabular}
\label{tab:bin_limit}
\end{table}
\end{savenotes}

\begin{table}[ht]
\centering
\caption{Formal\ae\ for the Binomial Model with Credit Risk}
\begin{tabular}{|rcl|l|l|}\hline
 \multicolumn{3}{|l|}{\textbf{Formul\ae}} & \textbf{Description} \\\hline
 \u     &=& $e^{\sigma\sqrt{\Dt}}$ &
    Multiplier for an up-step \\\hline
 \d     &=& $e^{-\sigma\sqrt{\Dt}}$ &
    Multiplier for a down-step \\\hline
 \pu    &=& $\frac{e^{r\Dt} - \d e^{-\lambda\Dt} - (1 - \eta)(1 - e^{-\lambda\Dt})}{\u - \d}$ &
    Probability of an up-step \\\hline
 \pd    &=& $e^{-\lambda\Dt} - \pu$ &
    Probability of a down-step \\\hline
 \po    &=& $1 - e^{-\lambda\Dt}$ &
    Probability of default \\\hline
 \V     &=& $e^{-r\Dt}(\Vu\pu + \Vd\pd + \X\po) + c_i$ &
    Value of portfolio at $t$ for $t^{c}_i = t$ \\\hline
\end{tabular}
\label{tab:bin_formulae}
\end{table}

\section{Example and Hedging Strategy}
Provide a numerical example (with diagrams) and provide example of the hedging strategy.


\chapter{Numerical Example}
Implement above models, explain implementation and results, by example.


\chapter{Conclusion}
Summarise what this paper shows and does (i.e. Introduction + Conclusion = Executive Summary).

% BIBLIOGRAPHY ---------------------------------------------------

\clearpage
\setlinespacing{1}
\addcontentsline{toc}{chapter}{Bibliography}
\bibliographystyle{../references/amsplnat}
\bibliography{../references/references.bib}
\end{document}

%
% Document ends here - everything beyond this point is ignored.
%


% GRAPHIC TEMPLATE FOR EPS OR PDF FILE ---------------------------

\begin{figure}[h!]
    \centering
    \includegraphics[width=8cm]{filename}
    \caption[Short caption for figure list]{Main caption.}\label{FigLabel}
\end{figure}

% GRAPHIC TEMPLATE FOR JPG FILE ----------------------------------

\begin{figure}[h!]
    \centering
    \includegraphics[\bb=0 0 XXX YYY, width=8cm]{filename.jpg}
    \caption[Short caption for figure list]{Main caption.}\label{FigLabel}
\end{figure}

% Replace XXX and YYY with pixel width and height of image.

